\documentclass[11pt,a4paper]{article}
\usepackage[T1]{fontenc}
\usepackage{lmodern,amsmath,amssymb,amsthm,mathtools}
\usepackage[margin=25mm]{geometry}
\usepackage{microtype,xcolor,listings,xurl,hyperref,fancyhdr}
\hypersetup{colorlinks=true,linkcolor=blue!45!black,citecolor=blue!45!black,urlcolor=blue!45!black,
pdftitle={An Elementary Proof of Sun's Weighted Prime-Square Congruence},
pdfauthor={Research manuscript prepared with ChatGPT for Keio}}
\setlength{\headheight}{14pt}
\pagestyle{fancy}\fancyhf{}
\fancyhead[L]{Sun's weighted prime-square congruence}
\fancyhead[R]{A finite telescoping proof}
\fancyfoot[C]{\thepage}
\newtheorem{theorem}{Theorem}[section]
\newtheorem{lemma}[theorem]{Lemma}
\newtheorem{corollary}[theorem]{Corollary}
\theoremstyle{remark}\newtheorem{remark}[theorem]{Remark}
\newcommand{\Z}{\mathbb Z}\newcommand{\Q}{\mathbb Q}
\lstset{language=Python,basicstyle=\ttfamily\footnotesize,
keywordstyle=\color{blue!55!black},commentstyle=\color{gray!80!black},
stringstyle=\color{gray!80!black},breaklines=true,columns=fullflexible,
keepspaces=true,showstringspaces=false,frame=single,
rulecolor=\color{gray!40},aboveskip=10pt,belowskip=10pt}
\title{\textbf{An Elementary Proof of Sun's\\Weighted Prime-Square Congruence}\\[5pt]
\large A reciprocal binomial sum and a reflected WZ certificate}
\author{Research manuscript prepared with ChatGPT for Keio}
\date{8 November 2026}
\begin{document}\maketitle

\begin{abstract}
Let \(a_n=\sum_k(-4)^k\binom nk^2\binom{n-k}k^2\).
We prove that, for every odd prime \(p\),
\(\sum_{n=0}^{p-1}(20n+17)a_n\equiv
p(7+10(\frac{-1}{p}))\pmod{p^2}\),
the second assertion of Sun's Conjecture 5.8(ii).
A universal congruence for sums of squared binomial coefficients
reduces the problem to a reciprocal Ramanujan-type finite sum.
Although its individual summands can have a factor of \(p\) in the
denominator, an explicit rational WZ certificate transforms the total
into integral boundary terms. The two surviving endpoints are
congruent to 3 and 4 modulo \(p\); every other term vanishes at the
required precision. The proof includes \(p=3\) uniformly and uses
no other supercongruence or all-integer divisibility theorem as an
input. We also deduce that \(p^2\) divides the original weighted sum
if and only if \(p\in\{3,17\}\). The established WZ antecedents are
credited, and historical priority is not asserted.
\end{abstract}

\section{The unchanged assertion}

Define
\begin{equation}\label{eq:original}
 a_n=\sum_{k=0}^{\lfloor n/2\rfloor}
       (-4)^k\binom nk^2\binom{n-k}k^2,
 \qquad
 S(N)=\sum_{n=0}^{N-1}(20n+17)a_n.
\end{equation}
The sequence is \(W_n(1,-4)\) in Sun's notation and OEIS A179524
\cite{Sun2010,OEIS}. The second assertion of Conjecture 5.8(ii) in
\cite[printed p.~33]{Sun2010} predicts
\begin{equation}\label{eq:target}
 S(p)\equiv p\left(7+10\left(\frac{-1}{p}\right)\right)
          \pmod{p^2}\qquad(p\text{ odd prime}).
\end{equation}
Here \((\frac{-1}{p})=(-1)^{(p-1)/2}\) is the Legendre symbol.
The sum in \eqref{eq:target} has exactly \(p\) terms.

The preceding research manuscript \cite{Previous} treated the first
assertion of Conjecture 5.8(ii), namely \(N\mid S(N)\) for every
positive integer \(N\). The present proof of \eqref{eq:target} is
independent of that manuscript. Its ingredients are finite binomial
identities, an explicit rational telescoping certificate, and
elementary prime-adic valuation estimates. In particular, every
ingredient needed for the new conclusion is proved below.

The certificate was found by reflecting familiar Ramanujan WZ
kernels. Such kernels and the WZ method have established antecedents
in Guillera's work \cite{Guillera}, Sun's proof of a Ramanujan
supercongruence \cite[Section~3]{SunRefinement}, and the parameter
pair reproduced by Hou, He, and Wang \cite[arXiv v4, (4.1)]{HHW}.
We give the resulting rational certificate directly, so its
validity does not require analytic continuation or convergence
of an infinite hypergeometric series.



\section{The companion congruence modulo a prime square}
\label{sec:companion}

In this section we prove the second assertion in Sun's
Conjecture 5.8(ii). The proof uses an explicit rational telescoping
certificate; it does not require the all-integer divisibility theorem
or a separate Ramanujan supercongruence as an input.

\begin{theorem}\label{p2:main}
For every odd prime \(p\),
\[
 \sum_{n=0}^{p-1}(20n+17)a_n
 \equiv p\left(7+10\left(\frac{-1}{p}\right)\right)
 \pmod{p^2}.
\]
\end{theorem}

We use congruences between rational numbers in the localization
\(\Z_{(p)}\). Every displayed rational congruence below is justified
by verifying the relevant denominators or by an explicit valuation.
Write \(b_m=\binom{2m}{m}\).

\subsection{A universal inner-sum congruence}

\begin{lemma}\label{p2:inner}
Let \(p\) be an odd prime and \(0\le m<p\). Then
\[
 \sum_{n=m}^{p-1}\binom nm^2
 \equiv \frac{p}{(2m+1)b_m}\pmod{p^2}.
\]
The rational number on the right belongs to \(\Z_{(p)}\), including
the cases in which its displayed denominator is divisible by \(p\).
\end{lemma}
\begin{proof}
The product linearization and the hockey-stick identity give
\[
 \binom nm^2=\sum_{j=0}^m
   \binom mj\binom{m+j}{j}\binom n{m+j},
\qquad
 \sum_{n=m}^{p-1}\binom nm^2
 =\sum_{j=0}^m d_j\binom p{m+j+1},
\]
where \(d_j=\binom mj\binom{m+j}{j}\).
The product linearization counts ordered pairs of \(m\)-subsets by
the size \(m+j\) of their union.
We compare the last sum termwise with
\[
             p\sum_{j=0}^m\frac{(-1)^{m+j}d_j}{m+j+1}.
\]
Put \(r=m+j+1\). If \(r<p\), use
\(\binom pr\equiv(-1)^{r-1}p/r\pmod{p^2}\).
If \(r=p\), the compared terms are exactly equal, since \(p\) is odd.
If \(r>p\), the binomial coefficient on the left vanishes,
whereas \(d_j\) is divisible by \(p\): indeed
\(p\le m+j\le2p-2\) and both \(m,j<p\). Also \(p<r<2p\), so the
compared term on the right is divisible by \(p^2\).

For completeness, the required finite rational identity is
\[
 \sum_{j=0}^m\frac{(-1)^{m+j}}{m+j+1}
                 \binom mj\binom{m+j}{j}
          =\frac1{(2m+1)b_m}.
\]
To prove it, expand
\[
 \frac1{m!}\frac{d^m}{dx^m}\bigl(x^m(1-x)^m\bigr)
       =\sum_{j=0}^m(-1)^j\binom mj\binom{m+j}{j}x^j.
\]
Multiply by \((-1)^mx^m\), integrate from zero to one, and integrate
by parts \(m\) times. The boundary terms vanish, and the result is
\(\int_0^1x^m(1-x)^m\,dx=m!^2/(2m+1)!\).

Finally, \(v_p((2m+1)b_m)\le1\). For \(m<(p-1)/2\) it is zero;
for \(m=(p-1)/2\) the factor \(2m+1=p\) contributes one; and for
\(m>(p-1)/2\) the central binomial coefficient contributes one while
\(2m+1\) is a unit. Thus the claimed rational congruence is well defined.
\end{proof}

\subsection{An exact rational telescoping certificate}

For nonnegative integers \(n,k\), define
\[
 c_n=\frac{(-4)^n b_n^2}{(4n+1)b_{2n}},\qquad
 \mathcal H(n,k)=c_n
       \frac{(-2n-\tfrac12)_k(n+\tfrac12)_k}{(\tfrac12)_k^2},
\]
where \((a)_k=a(a+1)\cdots(a+k-1)\) and \((a)_0=1\). Put
\[
 \mathcal D(n,k)=(20n+2k+7)\mathcal H(n,k),\qquad
 \mathcal E(n,k)=\frac{16(2n+1)^2}{4n-2k+3}\mathcal H(n,k),
\]
and set \(\mathcal E(-1,k)=0\).

\begin{lemma}[Rational certificate]\label{p2:wz}
For every \(n\ge0\) and every \(k\ge1\),
\[
 \mathcal D(n,k)-\mathcal D(n,k-1)
      =\mathcal E(n,k)-\mathcal E(n-1,k).
\]
Consequently, for every positive integer \(h\),
\begin{equation}\label{p2:rectangle}
 \sum_{n=0}^h(20n+7)c_n
       =\sum_{n=0}^h\mathcal D(n,h)
                     -\sum_{k=1}^h\mathcal E(h,k).
\end{equation}
\end{lemma}
\begin{proof}
For \(n\ge1\), successive-term ratios give
\[
 \frac{\mathcal H(n,k-1)}{\mathcal H(n,k)}
   =\frac{(2k-1)^2}{(2k-4n-3)(2n+2k-1)},
\]
\[
 \frac{\mathcal H(n-1,k)}{\mathcal H(n,k)}
 =-\frac{n(4n-2k-1)(4n-2k+1)}
          {8(2n-1)^2(2n+2k-1)}.
\]
After division by \(\mathcal H(n,k)\), the asserted certificate is
the rational identity
\[
\begin{split}
 20n+2k+7
 &-\frac{(20n+2k+5)(2k-1)^2}
          {(2k-4n-3)(2n+2k-1)}\\
 &=\frac{16(2n+1)^2}{4n-2k+3}
       +\frac{2n(4n-2k+1)}{2n+2k-1},
\end{split}
\]
which follows by clearing the two linear denominators.
For \(n=0\), use \(\mathcal H(0,0)=1\) and
\(\mathcal H(0,k)=-1/(2k-1)\) for \(k\ge1\) to check the same
identity directly with \(\mathcal E(-1,k)=0\).
Summing over \(0\le n\le h\) and \(1\le k\le h\) gives
\eqref{p2:rectangle}.
\end{proof}

\begin{theorem}[Reciprocal Ramanujan sum]\label{p2:reciprocal}
For every odd prime \(p\), with \(h=(p-1)/2\),
\[
 R_p:=\sum_{n=0}^h
       \frac{(20n+7)(-4)^n b_n^2}{(4n+1)b_{2n}}
       \in\Z_{(p)},\qquad R_p\equiv7\pmod p.
\]
\end{theorem}
\begin{proof}
We evaluate the right-hand side of \eqref{p2:rectangle} modulo \(p\).
For \(0\le n\le h\),
\[
 v_p(c_n)=
 \begin{cases}
 0,&4n+1<p,\\
 -1,&4n+1\ge p.
 \end{cases}
\]
Indeed \(b_n\) is a unit, and exactly one factor of \(p\) occurs
in \((4n+1)b_{2n}\) in the second range. It comes from \(4n+1\)
when \(4n+1=p\), and from \(b_{2n}\) otherwise.

For \(1\le n<h\), the product \((n+\tfrac12)_h\) contains
exactly one factor divisible by \(p\), namely \(p/2\).
If \(4n+1\ge p\), the other product
\((-2n-\tfrac12)_h\) also contains exactly one such factor,
namely \(-p/2\). If \(4n+1<p\), it contains none.
The denominator \((\tfrac12)_h^2\) is a unit. Therefore
\[
 v_p(\mathcal H(n,h))=1\quad(1\le n<h),
 \qquad \mathcal D(n,h)\in p\Z_{(p)}.
\]
This covers both ranges, including the possible equality \(4n+1=p\).

For \(1\le k\le h\), \(c_h\) has valuation \(-1\),
\((h+\tfrac12)_k=(p/2)_k\) has valuation one, and
\((-2h-\tfrac12)_k=(\tfrac12-p)_k\) and \((\tfrac12)_k\)
are units. Thus \(\mathcal H(h,k)\) is a unit. Since
\[
 4h-2k+3=2p+1-2k
\]
lies between \(p+2\) and \(2p-1\), it is a unit as well.
The factor \((2h+1)^2=p^2\) in \(\mathcal E(h,k)\) now gives
\[
                   \mathcal E(h,k)\in p^2\Z_{(p)}.
\]

Only the two endpoint terms remain. At \(n=0\),
\[
 \mathcal D(0,h)=-\frac{2h+7}{2h-1}
                   =-\frac{p+6}{p-2}\equiv3\pmod p.
\]
Cancellation of the finite products at \(n=h\) gives the exact identity
\[
 \mathcal H(h,h)=4^h b_h
                 \prod_{i=1}^{h-1}\frac{p+2i}{2i+1}.
\]
All denominators here are units. Also
\[
 \prod_{i=1}^{h-1}\frac{2i}{2i+1}
                        =\frac{4^h}{(2h)b_h}.
\]
Consequently, by Fermat's little theorem,
\[
 \mathcal H(h,h)\equiv\frac{16^h}{p-1}\equiv-1\pmod p,
\]
and hence
\[
 \mathcal D(h,h)=(22h+7)\mathcal H(h,h)
                 =(11p-4)\mathcal H(h,h)\equiv4\pmod p.
\]
Every term on the right-hand side of \eqref{p2:rectangle} is now
\(p\)-integral, and its reduction is \(3+4=7\).
The case \(p=3\) is included: the interior range is empty and
the displayed endpoint denominators are still units.
\end{proof}

\subsection{Returning to the original weighted sum}

\begin{proof}[Proof of Theorem~\ref{p2:main}]
Put \(\gamma_k=(-4)^k b_k^2\) and \(h=(p-1)/2\).
The elementary factorial identity
\(\binom n{2k}b_k=\binom nk\binom{n-k}k\) gives
\[
 a_n=\sum_k\gamma_k\binom n{2k}^2.
\]
For \(N>m\ge0\), direct telescoping gives
\begin{equation}\label{p2:weighted-inner}
 \sum_{n=m}^{N-1}(2n-m+1)\binom nm^2
                    =\frac{N^2}{m+1}\binom{N-1}m^2.
\end{equation}
Indeed, set \(f_m(n)=(n-m)^2\binom nm^2/(m+1)\) for \(n\ge m\).
Then \(f_m(n+1)=(n+1)^2\binom nm^2/(m+1)\), so the summand is
\(f_m(n+1)-f_m(n)\). The lower boundary \(f_m(m)\) is zero,
and the upper boundary has the displayed form.
Using \(20n+17=10(2n-2k+1)+(20k+7)\), interchange the finite sums
and apply \eqref{p2:weighted-inner} to obtain
\[
\begin{split}
 S(p)={}&10p^2\sum_{k=0}^h
          \frac{\gamma_k}{2k+1}\binom{p-1}{2k}^2\\
       &+\sum_{k=0}^h(20k+7)\gamma_k
                           \sum_{n=2k}^{p-1}\binom n{2k}^2.
\end{split}
\]
In the first sum, all terms with \(k<h\) vanish modulo \(p^2\).
The remaining term is \(10p\gamma_h\). Since
\[
 b_h=\binom{p-1}h\equiv(-1)^h\pmod p,\qquad
 (-4)^h\equiv(-1)^h\pmod p,
\]
we have \(\gamma_h\equiv(-1)^h=(\frac{-1}{p})\pmod p\).

Lemma~\ref{p2:inner}, with \(m=2k\), transforms the second sum modulo
\(p^2\) into
\[
 p\sum_{k=0}^h\frac{(20k+7)(-4)^k b_k^2}{(4k+1)b_{2k}}
                      =pR_p\equiv7p\pmod{p^2}
\]
by Theorem~\ref{p2:reciprocal}. Combining the two contributions proves
\[
                    S(p)\equiv p\left(7+10\left(\frac{-1}{p}\right)\right)
                    \pmod{p^2}.
\]
\end{proof}


\section{Consequences, scope, and verification}

\begin{corollary}\label{cor:exceptions}
For an odd prime \(p\),
\[
 p^2\mid S(p)\quad\Longleftrightarrow\quad p\in\{3,17\}.
\]
For every odd prime \(p\notin\{3,17\}\), one has
\(v_p(S(p))=1\).
\end{corollary}
\begin{proof}
If \(p\equiv1\pmod4\), Theorem~\ref{p2:main} gives
\(S(p)\equiv17p\pmod{p^2}\); this is zero precisely when \(p=17\).
If \(p\equiv3\pmod4\), it gives
\(S(p)\equiv-3p\pmod{p^2}\); this is zero precisely when \(p=3\).
For every other odd prime the residue has valuation exactly one.
\end{proof}

Together with the separate all-integer divisibility argument in
\cite{Previous}, the present theorem supplies proofs of both
assertions in Conjecture 5.8(ii). This statement concerns the
mathematical content of the two research manuscripts. No result
about the three unweighted cases in Conjecture 5.8(i) is claimed
here, and no assertion modulo \(p^3\) is inferred.

The essential denominator issue is explicit in
Theorem~\ref{p2:reciprocal}: some individual summands have valuation
\(-1\), so termwise reduction of the original reciprocal sum modulo
\(p\) would be invalid. The exact finite identity
\eqref{p2:rectangle} first replaces that sum by terms in
\(\Z_{(p)}\). The valuation estimates and the endpoint values then
give its residue. This is the arithmetic step beyond a change of
notation or a formal identity alone.

A bounded search of the exact sequence, the unchanged weight,
equivalent binomial notation, and the reciprocal sum did not identify
a previous complete proof of this assertion. That search does not
certify historical firstness. The WZ technique and the motivating
kernels are established mathematics and are credited above; we do
not claim to have invented that general method. This is a newly
prepared proof manuscript, checked internally by independent
readings and exact arithmetic. It has not undergone external peer
review or proof-assistant verification.

The program in the appendix audits the rational certificate, the
finite telescoping identity, all boundary valuations, the universal
inner-sum congruence, and the original weighted congruence. The
reported run passed 81 exact certificate checks and all stated
tests at the 25 odd primes at most 101. These finite computations
are not assumptions of any theorem above.

\appendix
\section{Exact reproducibility checks}

The following complete Python program uses only the standard
library. It is included inside this single TEX file. It checks
the proof's transcription using exact integers and rational
numbers; the unbounded result is supplied by the algebraic proof.
\begin{lstlisting}


"""Exact checks for the all-odd-prime companion congruence.
The proof is symbolic; these checks audit its transcription only.
"""
from fractions import Fraction as Q
from math import comb, prod, isqrt


def poch(a, k):
    return prod((a+j for j in range(k)), start=Q(1))


def b(n):
    return comb(2*n, n)


def c(n):
    return Q((-4)**n*b(n)**2, (4*n+1)*b(2*n))


def H(n, k):
    return c(n)*poch(Q(-4*n-1, 2), k)*poch(Q(2*n+1, 2), k)/poch(Q(1, 2), k)**2


def D(n, k):
    return (20*n+2*k+7)*H(n, k)


def E(n, k):
    if n == -1:
        return Q(0)
    return Q(16*(2*n+1)**2, 4*n-2*k+3)*H(n, k)


def vp(x, p):
    x = Q(x)
    if not x:
        return float('inf')
    ans = 0
    n, d = abs(x.numerator), x.denominator
    while n % p == 0:
        n //= p
        ans += 1
    while d % p == 0:
        d //= p
        ans -= 1
    return ans


for n in range(9):
    for k in range(1, 10):
        assert D(n, k)-D(n, k-1) == E(n, k)-E(n-1, k)

primes = [p for p in range(3, 102, 2)
          if all(p % q for q in range(2, isqrt(p)+1))]
partial = [0]
for n in range(101):
    a = sum((-4)**k*comb(n, k)**2*comb(n-k, k)**2
            for k in range(n//2+1))
    partial.append(partial[-1]+(20*n+17)*a)

for p in primes:
    h = (p-1)//2
    R = sum((20*n+7)*c(n) for n in range(h+1))
    assert R == sum(D(n, h) for n in range(h+1))-sum(E(h, k) for k in range(1, h+1))
    assert vp(R-7, p) >= 1
    assert D(0, h) == Q(-(p+6), p-2)
    assert H(h, h) == 4**h*b(h)*prod((Q(p+2*i, 2*i+1) for i in range(1, h)), start=Q(1))
    assert vp(H(h, h)+1, p) >= 1
    assert all(vp(H(n, h), p) == 1 for n in range(1, h))
    assert all(vp(E(h, k), p) == 2 for k in range(1, h+1))
    assert partial[p] % (p*p) == p*(7+10*(-1)**h) % (p*p)
    for m in range(p):
        T = sum(comb(n, m)**2 for n in range(m, p))
        assert vp(T-Q(p, (2*m+1)*b(m)), p) >= 2

print('PASS: 81 exact WZ identities, including n=0 boundary')
print('PASS: reciprocal congruence and all endpoint valuations at', len(primes), 'odd primes <= 101')
print('PASS: original companion congruence at the same primes')
print('PASS: every inner-sum congruence for m < p at those primes')


\end{lstlisting}

\begin{thebibliography}{9}\raggedright

\bibitem{Sun2010}
Zhi-Wei Sun,
\emph{On Ap\'ery numbers and generalized central trinomial coefficients},
arXiv:1006.2776v11, 19 August 2010, Conjecture 5.8(ii), printed p.~33.
\url{https://arxiv.org/abs/1006.2776v11}.
The specified version contains the conjecture; the subsequent
withdrawal notice concerns reorganization into other manuscripts.

\bibitem{OEIS}
OEIS Foundation Inc.,
\emph{The On-Line Encyclopedia of Integer Sequences}, entry A179524,
contributed by Zhi-Wei Sun, 17 July 2010; accessed 8 November 2026.
\url{https://oeis.org/A179524}.

\bibitem{Guillera}
Jes\'us Guillera,
\emph{Hypergeometric identities for 10 extended Ramanujan-type series},
Ramanujan Journal \textbf{15} (2008), 219--234.
\url{https://doi.org/10.1007/s11139-007-9074-0}.
Preprint: \url{https://arxiv.org/abs/1104.0396}.

\bibitem{SunRefinement}
Zhi-Wei Sun,
\emph{A refinement of a congruence result by van Hamme and Mortenson},
Illinois Journal of Mathematics \textbf{56} (2012), no.~3, 967--979.
\url{https://arxiv.org/abs/1011.1902}.
See Section~3 for the motivating WZ pair.

\bibitem{HHW}
Qinghu Hou, Haihong He, and Xiaoxia Wang,
\emph{Ramanujan-inspired series for \(1/\pi\) involving harmonic numbers},
Journal of Difference Equations and Applications \textbf{30} (2024),
775--786.
\url{https://doi.org/10.1080/10236198.2024.2325492}.
Preprint: \url{https://arxiv.org/abs/2305.00498v4}, equation~(4.1).

\bibitem{Previous}
Research manuscript prepared with ChatGPT for Keio,
\emph{A Proof of Sun's Divisibility Conjecture for A179524},
unpublished research manuscript, 8 November 2026, 13 pp.
The present paper does not use its theorem as an input.

\end{thebibliography}
\end{document}
